% ============================================================================
% SEC04.TEX
% Section 1.4: Lifecycle Pengembangan Game
% ============================================================================

\section{Lifecycle Pengembangan Game}

\begin{learningoutcomes}
\item Memahami fase-fase dalam lifecycle pengembangan game
\item Mengetahui aktivitas utama setiap fase
\item Memahami pentingnya planning dan documentation
\item Mampu mengidentifikasi milestone dalam development
\end{learningoutcomes}

\subsection{Pengertian Lifecycle Pengembangan Game}

\begin{definisi}[Game Development Lifecycle]
Lifecycle pengembangan game adalah proses sistematis dari konsep awal hingga rilis dan maintenance game, yang terdiri dari beberapa fase utama \cite{Gregory2018}.
\end{definisi}

\subsection{Fase-Fase Pengembangan Game}

\subsubsection{1. Pre-Production}

\textbf{Tujuan}: Membangun fondasi dan visi game

\textbf{Aktivitas Utama}:
\begin{itemize}
\item Brainstorming konsep
\item Market research
\item Prototyping core mechanics
\item Menyusun Game Design Document (GDD)
\item Budgeting dan resource planning
\item Team assembly
\end{itemize}

\textbf{Output}:
\begin{itemize}
\item GDD lengkap
\item Prototype yang dapat dimainkan
\item Art bible dan style guide
\item Technical design document
\item Project plan dan timeline
\end{itemize}

\subsubsection{2. Production}

\textbf{Tujuan}: Membangun konten dan fitur game

\textbf{Aktivitas Utama}:
\begin{itemize}
\item Asset creation (art, audio, models)
\item Programming dan implementation
\item Level design
\item Integration dan testing
\item Iterasi berdasarkan feedback
\end{itemize}

\textbf{Fase dalam Production}:
\begin{enumerate}
\item \textbf{Alpha}: Core features implemented, masih banyak bug
\item \textbf{Beta}: Feature-complete, focus pada bug fixing dan polish
\item \textbf{Gold Master}: Final version siap untuk rilis
\end{enumerate}

\subsubsection{3. Post-Production}

\textbf{Tujuan}: Rilis, marketing, dan maintenance

\textbf{Aktivitas Utama}:
\begin{itemize}
\item Final testing dan QA
\item Marketing dan promotion
\item Launch preparation
\item Distribution
\item Post-launch support
\item Updates dan patches
\item DLC development (jika ada)
\end{itemize}

\subsection{Metodologi Pengembangan}

\subsubsection{Waterfall Model}
Pendekatan linear dengan fase yang berurutan:
\begin{itemize}
\item \textbf{Keuntungan}: Terstruktur, mudah di-track
\item \textbf{Kekurangan}: Kurang fleksibel, sulit untuk perubahan
\end{itemize}

\subsubsection{Agile/Scrum}
Pendekatan iteratif dengan sprint pendek:
\begin{itemize}
\item \textbf{Keuntungan}: Fleksibel, adaptif terhadap perubahan
\item \textbf{Kekurangan}: Memerlukan komunikasi intensif
\end{itemize}

\subsubsection{Hybrid Approach}
Kombinasi waterfall untuk planning dan agile untuk execution.

\begin{catatan}
Kebanyakan studio game modern menggunakan pendekatan hybrid atau agile karena sifat kreatif dan iteratif dari pengembangan game.
\end{catatan}

\subsection{Milestone dan Deliverables}

\textbf{Milestone Penting}:
\begin{enumerate}
\item \textbf{Concept Approval}: Konsep disetujui untuk development
\item \textbf{Prototype Complete}: Prototype core mechanics selesai
\item \textbf{Pre-Alpha}: Basic gameplay loop functional
\item \textbf{Alpha}: Feature-complete, content masih kurang
\item \textbf{Beta}: Content-complete, focus pada polish
\item \textbf{Gold}: Siap untuk rilis
\item \textbf{Launch}: Game dirilis ke publik
\end{enumerate}

\subsection{Challenges dalam Development}

\begin{itemize}
\item \textbf{Scope Creep}: Fitur terus bertambah di tengah development
\item \textbf{Technical Challenges}: Masalah teknis yang tidak terduga
\item \textbf{Resource Constraints}: Keterbatasan waktu, budget, atau personel
\item \textbf{Communication Issues}: Masalah komunikasi dalam tim
\item \textbf{Market Changes}: Perubahan tren atau kompetitor
\end{itemize}

\subsection{Best Practices}

\begin{enumerate}
\item \textbf{Start Small}: Mulai dengan scope yang realistis
\item \textbf{Prototype Early}: Uji konsep sesegera mungkin
\item \textbf{Iterate Often}: Perbaiki berdasarkan feedback
\item \textbf{Document Everything}: Dokumentasi penting untuk konsistensi
\item \textbf{Test Continuously}: Testing harus dilakukan sepanjang development
\end{enumerate}
